\documentclass[11pt]{article}
\usepackage[margin=1in]{geometry}
\usepackage{amsmath,amssymb,amsthm}
\usepackage[hidelinks]{hyperref}
\usepackage{xurl}
\newtheorem{theorem}{Theorem}
\newtheorem{corollary}[theorem]{Corollary}
\newcommand{\barv}[1]{\overline{#1}}
\title{A Hermitian reflection obstruction for four alternating knots}
\author{Internal written proof and exact algebra controls}
\date{17 November 2026}
\begin{document}
\maketitle

The four knots
\[
 \mathcal S=\{13a_{373},13a_{375},13a_{391},13a_{392}\}
\]
have meridional rank four by their retained seven-sheeted cover
certificates. Their orthogonal reflection obstruction extends to
Hermitian reflections. This excludes sharp finite complex reflection
certificates whose meridians have order two, including targets outside
the real Coxeter classification. It asserts nothing about meridians
of order three or higher.

\section{Field, reflection and diagram conventions}

Let $k$ be a field of characteristic different from two, with an
involution $a\mapsto\bar a$, possibly trivial. Let $B$ be a Hermitian
form, conjugate-linear in the first variable and linear in the second.
It need not be nondegenerate or positive. For $B(u,u)=2$ set
\[
                         R_u(v)=v-B(u,v)u.
\]
This is an involution preserving $B$. Its displacement image is $ku$,
because $(R_u-I)u=-2u$. Thus for norm-two roots,
\[
 R_u=R_v\quad\Longleftrightarrow\quad
 v=\lambda u\text{ for some }\lambda\bar\lambda=1.
\]
Conjugate meridians have normals of one common nonzero norm, by
transporting a chosen normal through their conjugating isometries.
Rescaling $B$ by a scalar in its fixed field normalizes this norm to
two. No square root is required.

The literal PDs, arc numbering and original four-seed generating
traces are those of \path{UNIVERSAL_REFLECTION_OBSTRUCTION.tex}.
The present checker binds its four cases to
\path{seven_sheet_cover_inputs.json} and reconstructs the required
crossings. Named-census identification and the cover lower bounds
are separate retained dependencies.

\begin{theorem}\label{thm:hermitian}
For every $K\in\mathcal S$ and every representation of its knot group
by nonisotropic Hermitian reflections over $k$, the normal vectors
of its four specified generating meridians span a space of dimension
at most three. The form may be singular and the ambient dimension
arbitrary.
\end{theorem}

\section{Common normalization}

Suppose the original four normals are independent. For the three
siblings use the changed seed systems
\[
\begin{array}{c|c|c}
K&\text{original seeds}&\text{changed seeds}\\\hline
13a_{375}&(1,2,6,9)&(1,2,6,9)\\
13a_{391}&(1,2,5,9)&(1,2,6,9)\\
13a_{392}&(1,2,5,11)&(1,2,6,12).
\end{array}
\]
For $391$, putting $P=u_5,Q=u_9,z_0=B(P,Q)$ replaces the pair by
$(P-\bar z_0Q,Q)$. For $392$ the new pair is
\[
 u_6=(1-z_0\bar z_0)P+\bar z_0Q,
 \qquad u_{12}=Q-z_0P.
\]
Both changes have determinant one and consist of actual reflected
normals. Independence is preserved even if the form is singular.

Write the four new basis roots as $e_0=u_1,e_1=u_2,e_2=u_6,e_3$,
where $e_3=u_8$ for $373$, $u_9$ for $375,391$, and $u_{12}$ for
$392$. The common initial propagation is
\[
 (2,6,7),\ (7,2,3),\ (3,1,0),\ (1,3,4).
\]
With $t=B(e_1,e_2)$ this gives
\[
 u_3=(1-t\bar t)e_1+\bar t e_2,
 \quad u_0=e_0-su_3,\quad s=B(u_3,e_0).
\]
The unused crossing $(0,1,2)$ says
$R_{u_0}e_0=\varepsilon e_1$ with $\varepsilon\bar\varepsilon=1$.
Writing $H=B(u_0,e_0)$, comparison of the three independent
coefficients gives
\[
 H=1,\qquad s(1-t\bar t)=\varepsilon,\qquad s\bar t=0.
\]
Thus $s\ne0$, $t=0$ and $B(e_0,e_1)=\bar\varepsilon$.
Replace $e_1$ by $\varepsilon e_1$. Recomputing the trace normalizes
\[
 B(e_0,e_1)=1,\qquad u_0=U=e_0-e_1,
 \quad u_4=-U,\quad u_6=u_7=V=e_2.
\]
Put $W=e_3$. The roots $U,V,W$ are independent and have norm two.
Every unused relation below has a unit-norm phase in place of a
fixed sign; no phase is discarded.

\section{The four coefficient contradictions}

In all cases write $x=B(U,V)$, $y=B(U,W)$, $z=B(V,W)$.
For $373$ set $h=1-y\bar y$, $v=\bar x y+zh$ and
$w=B(U,u_{10})$. Its remaining trace gives
\[
\begin{aligned}
u_5&=-U+\bar yW,&u_9&=yU+hW,\\
u_{10}&=yU-vV+hW,&u_{11}&=(y-w)U-vV+hW,\\
u_{12}&=(1-\bar w y)U+\bar w vV-\bar w hW.
\end{aligned}
\]
The unused relation $R_{u_{11}}V=\beta W$, with $\beta\bar\beta=1$,
gives, on putting $p=B(u_{11},V)$,
\[
 p(y-w)=0,\qquad 1+pv=0,\qquad -ph=\beta.
\]
Hence $p\ne0$, $w=y$ and $v=\bar\beta h$.
The relation $R_{u_{12}}u_5=\alpha V$ now gives
$qh=-1$ and $\bar y\bar\beta=\alpha$, where $q=B(u_{12},u_5)$.
Since $\alpha\bar\alpha=\beta\bar\beta=1$, this forces
$y\bar y=1$, so $h=0$, contradicting $qh=-1$.

For $375$ put $T=R_VW=W-zV$, $t=B(U,T)$ and
$A=R_UT=T-tU$. The trace gives
\[
 u_8=R_TV,\quad u_{11}=A,\quad
 u_5=(t\bar t-1)U-\bar tT,\quad u_{12}=R_{u_8}A.
\]
As $u_8$ has no $U$-coefficient, $u_{12}$ has $U$-coefficient $-t$.
The unused crossing $R_Wu_{12}=\gamma U$ gives $t=-\gamma$.
Thus $t\bar t=1$ and $R_{u_5}=R_T$.
The other unused crossing would give
$\alpha W=R_{u_5}u_8=R_TR_TV=V$, a contradiction.

For $391$ keep $T,A,t$ and put $D=R_WV=V-\bar zW$.
The trace gives $u_5=D$, $u_{11}=A$, $u_{12}=R_DA$, and
$u_8=R_AV$. The unused relation says
\[
 u_{12}=\gamma R_TU=\gamma(U-\bar tT).
\]
Its $U$-coefficient gives $t=-\gamma$ and hence
$\gamma\bar t=-1$. The whole right side is therefore $A$.
Thus $u_{12}=A$ and the other unused crossing gives
$\beta W=R_Au_8=V$, a contradiction.

For $392$ put $h=1-z\bar z$, $t=B(U,u_{11})$ and
$k=B(V,u_{10})$. The trace gives
\[
\begin{aligned}
u_5&=V-\bar zW,&u_{11}&=zV+hW,\\
u_{10}&=-tU+zV+hW,&u_9&=-tU+(z-k)V+hW,\\
u_8&=\bar k tU+(1-\bar k z)V-\bar k hW.
\end{aligned}
\]
The relation $R_{u_9}W=\gamma U$ yields
$pt=\gamma$, $p(z-k)=0$ and $ph=1$, so $p,t,h\ne0$ and $k=z$.
The $U$-coefficient in $R_{u_{11}}u_8=\beta u_9$ gives
$\bar k t=-\beta t$. Thus $\bar k=-\beta$ has unit norm,
forcing $h=1-z\bar z=0$, again a contradiction.
This proves Theorem~\ref{thm:hermitian}.

\section{Finite complex images and exact assurance}

\begin{corollary}\label{cor:finite}
Let $K\in\mathcal S$. If a finite-image complex representation
$\rho$ sends meridians to order-two rank-one pseudo-reflections,
then
\[
 D=\sum_{g\in\rho(\pi_1)}\operatorname{im}(g-I)
 \quad\text{has dimension at most three}.
\]
In particular such a representation on a four-dimensional space
without a nonzero common fixed vector cannot exist. No finite
complex reflection quotient of essential dimension four can give
the sharp meridional bound for these knots.
\end{corollary}
\begin{proof}
Average any positive definite Hermitian form over the finite image.
Each meridian is then a unitary involution with a one-dimensional
minus-one eigenspace, hence a Hermitian reflection. Theorem~\ref{thm:hermitian}
gives a space $U$ of dimension at most three spanned by the four
generating normals. Every generator preserves $U$ and displaces
every vector into $U$; induction gives $D\subset U$.
Conversely each normal belongs to $D$ because its reflection
displaces it by minus twice itself. For a unitary representation,
$D^\perp$ is exactly the common fixed space, proving the last claim.
The averaging proof needs no classification of complex reflection
groups. An order-three pseudo-reflection is outside these hypotheses.
\end{proof}

The checker \path{verify_hermitian_reflection_obstruction.py}
uses exact sparse Laurent polynomials with independent conjugate
pairs and unit phases satisfying $\bar\alpha=\alpha^{-1}$,
$\bar\beta=\beta^{-1}$, $\bar\gamma=\gamma^{-1}$.
It reconstructs the literal four trace systems, checks the root
formulas and norm identities, and verifies a constant-one identity
in the conjugation-stable ideal for $373$. This covers arbitrary
unit phases, not a finite sample. The remaining coefficient
cancellations are the written arguments above.
The input identities and standard finite averaging argument are
internal mathematical proofs with exact algebra controls, not
formal verification or external specialist review. Worldwide priority
and noninvolutory lower certificates remain separate questions.
\end{document}
